Au lycée, les mathématiques changent d'objet : on quitte le nombre pour la fonction, et le
calcul exact pour la limite. Les dix chapitres suivent le programme de la filière S — le
dernier avant la réforme de 2019, dont le contenu reste presque intégralement en place dans
la spécialité mathématiques actuelle.

C'est ici que la frontière de l'\emph{admis} devient le sujet principal. Le programme du
lycée admet le théorème des valeurs intermédiaires, l'existence des primitives, la formule
de l'intégrale comme aire, le théorème de Moivre--Laplace. Un assistant de preuve ne connaît
pas cette convention : chacun de ces résultats doit être démontré, ou emprunté à une
bibliothèque qui l'a fait. Le déplacement se lit chapitre par chapitre, et il va dans les
deux sens — l'espérance d'une loi binomiale, admise en première, est démontrée ici, tandis
que le théorème du toit, admis en terminale, ne l'est pas.

L'analyse occupe la moitié de la partie, et sa difficulté n'est pas où on l'attend. Les
opérations sur les limites sont faciles ; ce qui coûte, c'est de vérifier que la définition
avec $\varepsilon$ du cours et celle de la bibliothèque, écrite avec des filtres, sont bien
la même. Les premiers énoncés de plusieurs chapitres ne font que cela.
